\subsection{MVDR 的闭式解及最大感知 SINR 的证明}

\noindent\textbf{命题。}任取 UAV $u_i$ 和时隙 $t$，固定 UAV 位置、目标、频谱共享关系、CU 发射功率、相关信道及感知发射波束。设 $\sigma^2>0$、$p_j(t)\ge0$、$\bm a_r(\theta_i(t))\ne\bm0$，且感知信道满足式~\eqref{eq:A}。那么，式~\eqref{eq:mvdr} 所示单位范数接收波束能够达到该时隙最大的感知 SINR。其方向与经典无失真 MVDR 接收波束相同。

\medskip
\noindent\textbf{证明。}为简化符号，在本命题的证明中省略下标 $i$ 与时隙 $t$，记接收导向矢量为 $\bm a=\bm a_r(\theta_i(t))$，感知发射波束为 $\bm w$，并定义干扰与噪声协方差矩阵
\begin{equation}\label{eq:deriv-R}
\mathbf R=\sum_{j=1}^J\eta_{i,j}(t)p_j(t)
\bm h_{c_j,u_i}^{H}(t)\bm h_{c_j,u_i}(t)+\sigma^2\mathbf I.
\end{equation}
对于任意 $\bm v\ne\bm0$，由 $\eta_{i,j}(t)\ge0$、$p_j(t)\ge0$ 和 $\sigma^2>0$，有
\begin{equation}\label{eq:deriv-positive}
\bm v^H\mathbf R\bm v
=\sum_{j=1}^J\eta_{i,j}(t)p_j(t)
\left|\bm h_{c_j,u_i}(t)\bm v\right|^2+\sigma^2\|\bm v\|^2>0.
\end{equation}
因此 $\mathbf R$ 为 Hermitian 正定矩阵，其逆矩阵及正定平方根均存在。

首先，由式~\eqref{eq:A} 的秩一感知信道表达式，并利用接收波束约束 $\|\bm g\|=1$，感知 SINR 可以写为
\begin{equation}\label{eq:deriv-sinr}
\operatorname{SINR}(\bm g,\bm w)
= K(\bm w)\frac{|\bm g^H\bm a|^2}{\bm g^H\mathbf R\bm g},
\quad
K(\bm w)=\xi_0\rho d_i^{-4}(t)
\left|\bm a_t^H(\theta_i(t))\bm w\right|^2\ge0.
\end{equation}
这里 $\sigma^2=\sigma^2\bm g^H\bm g$ 使用了单位范数条件；后文将分式视为对非零向量的齐次扩展。由于 $K(\bm w)$ 与 $\bm g$ 无关，若 $K(\bm w)>0$，最大化感知 SINR 等价于求解
\begin{equation}\label{eq:deriv-max}
\max_{\|\bm g\|=1}
\frac{|\bm g^H\bm a|^2}{\bm g^H\mathbf R\bm g}.
\end{equation}
由复数 Cauchy--Schwarz 不等式，对任意非零 $\bm g$ 有
\begin{align}\label{eq:deriv-CS}
|\bm g^H\bm a|^2
&=\left|\left(\mathbf R^{1/2}\bm g\right)^H
\left(\mathbf R^{-1/2}\bm a\right)\right|^2\notag\\
&\le \left(\bm g^H\mathbf R\bm g\right)
\left(\bm a^H\mathbf R^{-1}\bm a\right).
\end{align}
将两侧除以正数 $\bm g^H\mathbf R\bm g$ 得到最大值上界。式~\eqref{eq:deriv-CS} 取等当且仅当 $\bm g=\beta\mathbf R^{-1}\bm a$，其中 $\beta\ne0$。故单位范数约束下，所有最大 SINR 解可写为
\begin{equation}\label{eq:deriv-unit}
\bm g^{\max}=e^{\mathrm j\phi}\,
\frac{\mathbf R^{-1}\bm a}{\|\mathbf R^{-1}\bm a\|},
\qquad \phi\in\mathbb R,
\end{equation}
其最大 SINR 为 $K(\bm w)\bm a^H\mathbf R^{-1}\bm a$。选择 $\phi=0$ 并恢复下标，即得到式~\eqref{eq:mvdr}。因为式~\eqref{eq:A} 具有秩一结构，$\bm w$ 只出现在标量 $K(\bm w)$ 中，所以最大 SINR 的接收方向与 $\bm w$ 的取值无关。若 $K(\bm w)=0$，所有单位范数接收波束的感知 SINR 均为零，式~\eqref{eq:deriv-unit} 仍是最大 SINR 的一个解；若感知阈值 $\epsilon>0$，此时原方案不可行。

接下来证明上述方向与经典无失真 MVDR 相同。经典 MVDR 求解问题为
\begin{equation}\label{eq:deriv-mvdr-problem}
\min_{\bm v\in\mathbb C^N}\ \bm v^H\mathbf R\bm v
\quad\text{s.t.}\quad \bm v^H\bm a=1.
\end{equation}
该问题在保持目标方向单位响应的条件下，最小化输出干扰与噪声功率；参见 Capon [1] 及 Vorobyov [2，第 2.2 节]。根据式~\eqref{eq:deriv-CS}，任意满足约束的 $\bm v$ 均有
\begin{equation}\label{eq:deriv-mvdr-bound}
1=|\bm v^H\bm a|^2
\le (\bm v^H\mathbf R\bm v)(\bm a^H\mathbf R^{-1}\bm a),
\qquad
\bm v^H\mathbf R\bm v\ge
\frac{1}{\bm a^H\mathbf R^{-1}\bm a}.
\end{equation}
其中 $\bm a^H\mathbf R^{-1}\bm a>0$。令
\begin{equation}\label{eq:deriv-classic}
\bm v^{\rm MVDR}=
\frac{\mathbf R^{-1}\bm a}{\bm a^H\mathbf R^{-1}\bm a},
\end{equation}
则 $(\bm v^{\rm MVDR})^H\bm a=1$，并且式~\eqref{eq:deriv-mvdr-bound} 取等，因此式~\eqref{eq:deriv-classic} 是经典 MVDR 问题的最优解。式~\eqref{eq:deriv-classic} 与式~\eqref{eq:deriv-unit} 的方向相同：前者采用无失真归一化，后者采用论文要求的单位范数归一化，两者的尺度约束不同。

实际计算时，可先求解线性方程组 $\mathbf R\bm x=\bm a$，再根据需要取 $\bm v^{\rm MVDR}=\bm x/(\bm a^H\bm x)$ 或 $\bm g^{\max}=\bm x/\|\bm x\|$，无需显式计算 $\mathbf R^{-1}$。上述结论适用于当前给定的信道、匹配、功率和目标方向；当这些量变化时，需相应更新 $\mathbf R$ 和 $\bm a$。命题得证。\hfill$\square$

